\section{Research Questions and Hypotheses}\label{research-questions-and-hypotheses}

The study addresses the following primary research question:

\begin{quote}
\textbf{How does TabPFN compare with established gradient-boosted tree algorithms across different training-data regimes in tabular classification?}
\end{quote}

To address the primary question, four supporting research questions are examined:

\begin{enumerate}
\def\labelenumi{\arabic{enumi}.}
\item
  \textbf{Does TabPFN provide an advantage when training data are limited?}
\item
  \textbf{How does the relative performance of TabPFN change as training size increases?}
\item
  \textbf{Is the performance advantage of TabPFN consistent across datasets and evaluation metrics?}
\item
  \textbf{What computational trade-offs exist between TabPFN and gradient-boosted tree models?}
\end{enumerate}

Based on these research questions, four hypotheses were investigated:

\subsubsection{H1}\label{h1}

\textbf{TabPFN will be competitive or superior to gradient-boosted tree models when training data are limited.}

\subsubsection{H2}\label{h2}

\textbf{TabPFN's relative advantage will change as training size increases.}

\subsubsection{H3}\label{h3}

\textbf{TabPFN's performance relative to tree-based models will vary across datasets and evaluation metrics.}

\subsubsection{H4}\label{h4}

\textbf{Predictive advantages will be accompanied by computational trade-offs.}
